% !TeX root = ../main.tex

\chapter{章名（章節示例）}
章內容內容內容內容內容 \\
內容內容內容

\section{節名}
節內容內容內容內容內容 \\
內容內容內容

\subsection{小節名}
內容內容內容 \\
內容內容內容

\subsubsection{小小節}
內容內容內容 \\
內容內容內容

\paragraph{段}
內容內容內容 \\
內容內容內容

\subparagraph{小段}
內容內容內容 \\
內容內容內容


\section{文字}
第一行。
仍是第一行。 \\
第二行。


\section{圖片}

\subsection{插入圖片}

\begin{enumerate}
    \item 附錄中的圖應編為「附圖」，並須編入圖目次中。    
    \item 附錄中的圖編碼方式應與正文中的圖編碼一致。例如：
    \begin{itemize}
        \item 正文中「圖一」，附錄則用「附圖一」；
        \item 正文中「圖 1」，附錄則用「附圖 1」，以此類推。
    \end{itemize}    
    \item 附圖標題置於正下方，並且中對齊。    
    \item 如引用他人資料，務必於附圖標題下方註明來源，並且置中對齊。
\end{enumerate}

\begin{figure}[htbp]
    \centering
    \includegraphics[width=0.2\textwidth]{seal.eps}
    \caption{插入單一圖片width0.2}
    \label{fig:insert-single-image}
\end{figure}

\subsection{插入單一圖片}
可透過調整width來設定圖片的寬度

\begin{figure}[htbp]
    \centering
    \includegraphics[width=0.4\textwidth]{seal.eps}
    \caption{插入單一圖片width0.4}
    \label{fig:insert-single-image}
\end{figure}


\newpage


\subsection{插入多張圖片}
引用網路上的圖片，根據學術引用規範（例如 APA、MLA、Chicago 等），應該標註圖片來源。你可以在圖片的標題或參考文獻中加入來源, 或以footnotetext加注於頁下 ：

\begin{figure}[ht]
\centering
\includegraphics[width=0.48\linewidth]{seal.eps}
\hfill
\includegraphics[width=0.48\linewidth]{seal.png}
\caption{插入多張圖片 \protect\footnotemark}
\end{figure}

\footnotetext{圖片來源：\href{https://waxinvest.com/projects/abundant-robots/}{Abundant Robotics 官方網站}（訪問日期：2025/1/16）。}


